\begin{claim}{AC.39}{Proved: exact finite-level meaning of hereditary relative tests}
Let \(G\) be a finite solvable group of order prime to \(p\),
and let \(V\) be a finite \(p\)-primary abelian group with a
\(G\)-action. For a subgroup \(H\), write
\[
                 e_H=\frac1{|H|}\sum_{h\in H}h
                      \quad\text{on }V.
\]
The following conditions are equivalent:
\begin{enumerate}
\item For every pair of normal subgroups \(H<K\) of \(G\),
      one has \(V^H\supsetneq V^K\).
\item There is \(v\in V\) such that
      \(\operatorname{Stab}_G(e_Hv)=H\) for every normal
      subgroup \(H\) of \(G\).
\end{enumerate}

In particular, let \(D/k\to C/k\) be a constant solvable
prime-to-\(p\) Galois etale cover, with a rational chosen
point over \(P\), and put \(V=J_D(k)[p^\infty]\).
All relative class-number tests for
\[
                    D/H\longrightarrow D/K,
                    \qquad H<K\text{ normal in }G,
\]
fail if and only if there is a rational \(p\)-primary
class on \(D\) whose norms have free deck orbits on
every normal quotient of this particular cover.
This is a finite-level statement about classes, with
no assertion that they are differences of points.
\end{claim}
\begin{proof}
Necessity is immediate: if \(V^H=V^K\), then \(e_Hv\)
is \(K\)-fixed, so cannot have stabilizer exactly \(H\).
For sufficiency, induct on \(|G|\). The trivial group
causes no difficulty. Choose a nontrivial abelian normal
subgroup \(N\) of the solvable group \(G\), and split
\[
                  V=U\oplus W,\qquad
                  U=e_NV=V^N,\quad W=(1-e_N)V.
\]
Both summands are \(G\)-stable. The strict-invariant
condition descends to the \(G/N\)-module \(U\), so induction
gives \(u\in U\) with the required property for all normal
quotients of \(G/N\). For any normal \(H\) of \(G\),
\(e_Hu=e_{HN}u\), whence
\[
                     \operatorname{Stab}_G(e_Hu)=HN.
\]

We choose \(w\in W\) simultaneously for all the normal
subgroups \(H\). The algebra \(\mathbf Z_p[N]\) is finite
etale and commutative; it is a product of unramified
discrete valuation rings, and \(N\) acts through a
prime-to-\(p\) character on each factor. The projectors
\(e_H\) are central in \(\mathbf Z_p[G]\). They commute
with this algebra and with each other, so decompose
\(W\) further into their finitely many joint image/kernel
summands. Choose one nonzero element in every nonzero
joint summand and add these elements to obtain \(w\).

If a prime-to-\(p\) root of unity is not one, its
difference from one is a unit in the relevant
unramified valuation ring. Thus every nonzero element
of a character summand has precisely that character's
kernel as its stabilizer in \(N\), even though the
module need not be free over that ring. Our choice
therefore ensures, simultaneously for all normal \(H\),
\[
 \operatorname{Stab}_N(e_Hw)
       =\ker\bigl(N\longrightarrow
                         \operatorname{Aut}(W^H)\bigr).
\]
Denote this kernel by \(K_H\). It is normal in \(G\),
because \(N\) and \(H\) are normal and \(W^H\) is
\(G\)-stable. It contains \(H\cap N\). If this
containment were strict, then \(HK_H>H\); both would
be normal in \(G\), while \(K_H\) acts trivially on
both \(U^H\) and \(W^H\). This would give
\(V^H=V^{HK_H}\), contrary to the hypothesis.
Consequently \(K_H=H\cap N\).

Set \(v=u+w\). Since \(U\) and \(W\) are \(G\)-stable,
an element stabilizing \(e_Hv\) must stabilize each
component separately. From the first component it
belongs to \(HN\). Modulo \(H\) it is therefore
represented by some \(n\in N\); from the second
component this \(n\) lies in \(H\cap N\).
The stabilizer is exactly \(H\), as required.

For the geometric application, the prime-to-\(p\)
pullback and norm identities identify
\(J_{D/H}(k)[p^\infty]\) with \(V^H\). Therefore
\[
 v_p\!\left(\frac{h(D/H)}{h(D/K)}\right)>0
                  \quad\Longleftrightarrow\quad
                  V^H\supsetneq V^K.
\]
The pullback of the actual norm of a class \(v\)
is \(|H|e_Hv\). Multiplication by \(|H|\) is
invertible on \(V\), so it does not change its
stabilizer. The abstract equivalence gives exactly
the asserted finite-level application.
\end{proof}

\paragraph{Why tests only over the original base are weaker.}
Let \(p=2\), \(G=C_7\rtimes C_3\), and let \(V\) be
the additive group of \(\mathbf F_4\), with the action
inflated from \(G^{\rm ab}=C_3\) by multiplication by
a primitive cube root. The proper normal subgroups
are \(1\) and \(C_7\), and
\[
                     V^G=0,\qquad
                     V^1=V^{C_7}=V\ne0.
\]
Thus every nontrivial normal quotient over the original
base has a nonzero new module, and all norm maps have
their usual prime-to-two splittings. Nevertheless every
vector is fixed by \(C_7\), so no free orbit exists.
The omitted relative test is exposed by
\(V^1=V^{C_7}\). This is an algebraic countermodel to
the inference from those module conditions alone;
it is not asserted to be the full Jacobian module of
a specified curve or a counterexample to A1 or A10.

\paragraph{Remaining inverse-system issue.}
AC.39 treats a finite cover with constant deck group.
At a fixed finite field, such covers need not be
cofinal in the geometric prime-to-\(p\) tower.
For nonconstant deck groups, the rational subgroup
\(J_D(k)[p^\infty]\) need not be \(G\)-stable, so the
commuting-projector argument cannot simply be repeated
on that subgroup. Making the deck group constant
produces regular classes over an extension, but does
not make them fixed by the original Frobenius.
Moreover regularity on a cofinal selection of levels
does not imply regularity after norms to every
geometric normal quotient. A semilinear descent and
compatibility statement resolving these two issues
is still required before AC.39 could imply a
Frobenius-fixed element of the full ghost module.
